\begin{helper}
Let $d\sigma$ denote Lebesgue measure on $\mathbb{C}\cong\mathbb{R}^{2}$ and
let $g\in L^{1}(\mathbb{C})$, i.e. $g\colon\mathbb{C}\to\mathbb{C}$ is
integrable for $d\sigma$.  For each $m\in\mathbb{N}$ let $\varphi_m$ be a
smooth bump function centred at $0$: a $C^\infty$ function
$\varphi_m\colon\mathbb{C}\to\mathbb{R}$ with radii
$0<r_m^{\mathrm{in}}<r_m^{\mathrm{out}}$ such that $0\le\varphi_m\le1$,
$\varphi_m\equiv1$ on $\{|y|\le r_m^{\mathrm{in}}\}$ and $\varphi_m(y)=0$ for
$|y|\ge r_m^{\mathrm{out}}$.  Write
\[
  \varrho_m=\frac{\varphi_m}{\int\varphi_m\,d\sigma}
\]
for its normalisation, so that $\varrho_m\in C^\infty(\mathbb{C})$,
$\varrho_m\ge0$, $\varrho_m(y)=0$ for $|y|\ge r_m^{\mathrm{out}}$, and
$\int\varrho_m\,d\sigma=1$.  Assume
\[
  r_m^{\mathrm{out}}\le1\quad\text{for every }m,
  \qquad\text{and}\qquad r_m^{\mathrm{out}}\longrightarrow0 .
\]
Then, with $(\varrho_m*g)(x)=\int\varrho_m(t)\,g(x-t)\,d\sigma(t)$,
\[
  \int_{\mathbb{C}}\bigl|(\varrho_m*g)(x)-g(x)\bigr|\,d\sigma(x)
  \longrightarrow0\qquad(m\to\infty).
\]

This is the $L^{1}$-norm form of the convergence of mollifications.  Only the
pointwise form (at a point of continuity) and the almost-everywhere form are
standard background; the $L^{1}$ form is what is needed whenever the limit
$m\to\infty$ has to be taken inside a further integration, for there
almost-everywhere convergence of the $\varrho_m*g$ would have to be
accompanied by an $m$-uniform dominating function, which mere local
integrability of $g$ does not supply.  The normalisation
$\int\varrho_m\,d\sigma=1$ is what makes the statement true, and the bound
$r_m^{\mathrm{out}}\le1$ only serves to confine all the supports involved to
one fixed ball.
\end{helper}
\begin{proof}
Throughout, all integrals without an indicated domain are over $\mathbb{C}$
against $d\sigma$.  Note first that each $\varrho_m$ is well defined, since
$\varphi_m\ge0$ is continuous, is $\ge1$ on the ball $\{|y|\le
r_m^{\mathrm{in}}\}$ of positive measure and vanishes off a bounded set, so
$0<\int\varphi_m\,d\sigma<\infty$; and that $\varrho_m\in L^{1}(\mathbb{C})$
with $\int|\varrho_m|\,d\sigma=\int\varrho_m\,d\sigma=1$, because
$\varrho_m\ge0$.

\emph{Step 1: convolution with an $L^{1}$ kernel is an $L^{1}$ contraction
(Young's inequality at $p=1$).}  We claim that for every real-valued
$\rho\in L^{1}(\mathbb{C})$ and every $f\in L^{1}(\mathbb{C})$,
\begin{equation}\tag{1}
  \int\bigl|(\rho*f)(x)\bigr|\,d\sigma(x)
  \le\Bigl(\int|\rho(t)|\,d\sigma(t)\Bigr)
     \Bigl(\int|f(x)|\,d\sigma(x)\Bigr).
\end{equation}
Indeed, the map $(x,t)\mapsto\rho(t)f(x-t)$ is measurable on
$\mathbb{C}\times\mathbb{C}$, being the product of a measurable function of
$t$ with the composition of the measurable $f$ with the continuous map
$(x,t)\mapsto x-t$; and it is integrable for $d\sigma\otimes d\sigma$, since
by Tonelli's theorem and the translation invariance of $d\sigma$
\[
  \iint|\rho(t)|\,|f(x-t)|\,d\sigma(x)\,d\sigma(t)
  =\int|\rho(t)|\Bigl(\int|f(x-t)|\,d\sigma(x)\Bigr)d\sigma(t)
  =\Bigl(\int|\rho|\Bigr)\Bigl(\int|f|\Bigr)<\infty .
\]
In particular $(\rho*f)(x)=\int\rho(t)f(x-t)\,d\sigma(t)$ is an absolutely
convergent integral for almost every $x$, the function $\rho*f$ is in
$L^{1}(\mathbb{C})$, and
\[
  \int\bigl|(\rho*f)(x)\bigr|\,d\sigma(x)
  \le\int\Bigl(\int|\rho(t)|\,|f(x-t)|\,d\sigma(t)\Bigr)d\sigma(x)
  =\Bigl(\int|\rho|\Bigr)\Bigl(\int|f|\Bigr),
\]
where the inequality is the triangle inequality for the inner integral
and the equality is the computation just displayed, read with the order of
integration exchanged by Fubini's theorem.  This proves (1).  Applied to
$\rho=\varrho_m$, for which $\int|\varrho_m|=1$, it gives
\begin{equation}\tag{2}
  \int\bigl|(\varrho_m*f)(x)\bigr|\,d\sigma(x)\le\int|f|\,d\sigma
  \qquad\bigl(f\in L^{1}(\mathbb{C}),\ m\in\mathbb{N}\bigr).
\end{equation}

\emph{Step 2: the assertion for a continuous compactly supported function.}
Let $h\colon\mathbb{C}\to\mathbb{C}$ be continuous with compact support.  We
claim
\begin{equation}\tag{3}
  \int\bigl|(\varrho_m*h)(x)-h(x)\bigr|\,d\sigma(x)\longrightarrow0 .
\end{equation}
Being continuous with compact support, $h$ is bounded, say $|h|\le C$, and
integrable; choose $R>0$ with $\operatorname{supp}h\subseteq\{|y|\le R\}$ and
put
\[
  S=\{x\in\mathbb{C}:|x|\le1+R\},
\]
a closed ball, hence measurable with $\sigma(S)<\infty$.

First, both $h$ and $\varrho_m*h$ vanish off $S$.  For $h$ this is clear from
$R\le1+R$.  For $\varrho_m*h$, the support of a convolution is contained in
the closure of the sum of the supports, so if $(\varrho_m*h)(x)\ne0$ then
$x=a+b$ with $|a|<r_m^{\mathrm{out}}$ and $|b|\le R$, whence
$|x|\le|a|+|b|<r_m^{\mathrm{out}}+R\le1+R$ by the hypothesis
$r_m^{\mathrm{out}}\le1$, i.e. $x\in S$.

Second, on all of $\mathbb{C}$ we have the $m$-independent bound
\begin{equation}\tag{4}
  \bigl|(\varrho_m*h)(x)-h(x)\bigr|\le2C+1 .
\end{equation}
Indeed, since $\int\varrho_m\,d\sigma=1$,
\[
  (\varrho_m*h)(x)-h(x)=\int\varrho_m(t)\bigl[h(x-t)-h(x)\bigr]\,d\sigma(t),
\]
and $\varrho_m\ge0$ vanishes for $|t|\ge r_m^{\mathrm{out}}$, so the modulus
of the left-hand side is at most
$\sup_{|y-x|<r_m^{\mathrm{out}}}|h(y)-h(x)|\le2C\le2C+1$.

Third, for each fixed $x$ we have $(\varrho_m*h)(x)\to h(x)$.  This is the
pointwise form of the convergence of mollifications: given $\epsilon>0$,
continuity of $h$ at $x$ provides $\tau>0$ with $|h(y)-h(x)|\le\epsilon$
whenever $|y-x|\le\tau$; since $r_m^{\mathrm{out}}\to0$ we have
$r_m^{\mathrm{out}}\le\tau$ for all large $m$, and then the display above
bounds $|(\varrho_m*h)(x)-h(x)|$ by
$\int\varrho_m(t)\,\epsilon\,d\sigma(t)=\epsilon$.

Now apply the dominated convergence theorem to the functions
$x\mapsto|(\varrho_m*h)(x)-h(x)|$, which are measurable (each
$\varrho_m*h$ is in $L^{1}$ by Step 1, and $h$ is measurable).  By the three
observations they converge pointwise to $0$ and are dominated by the fixed
integrable function $(2C+1)\mathbf1_S$: on $S$ this is (4), and off $S$ the
function is identically $0$.  Hence their integrals tend to $0$, which is (3).

\emph{Step 3: the general case, by density.}  Let $\eta>0$.  The continuous
compactly supported functions are dense in $L^{1}(\mathbb{C})$, so there is a
continuous $h$ with compact support such that
\begin{equation}\tag{5}
  \int|g-h|\,d\sigma\le\frac{\eta}{3};
\end{equation}
such an $h$ is in particular integrable.  Fix $m$.  Since
$\varrho_m\in L^{1}$ and $g,h\in L^{1}$, Step 1 shows that for almost every
$x$ both defining integrals of $(\varrho_m*g)(x)$ and $(\varrho_m*h)(x)$
converge absolutely, so for almost every $x$
\[
  (\varrho_m*g)(x)-(\varrho_m*h)(x)
  =\int\varrho_m(t)\bigl[g(x-t)-h(x-t)\bigr]\,d\sigma(t)
  =\bigl(\varrho_m*(g-h)\bigr)(x).
\]
Combining this with (2) applied to $f=g-h\in L^{1}$ and with (5),
\begin{equation}\tag{6}
  \int\bigl|(\varrho_m*g)(x)-(\varrho_m*h)(x)\bigr|\,d\sigma(x)
  =\int\bigl|\bigl(\varrho_m*(g-h)\bigr)(x)\bigr|\,d\sigma(x)
  \le\int|g-h|\,d\sigma\le\frac{\eta}{3}.
\end{equation}
The pointwise triangle inequality
\[
  \bigl|(\varrho_m*g)(x)-g(x)\bigr|
  \le\bigl|(\varrho_m*g)(x)-(\varrho_m*h)(x)\bigr|
    +\bigl|(\varrho_m*h)(x)-h(x)\bigr|
    +\bigl|h(x)-g(x)\bigr|
\]
holds for every $x$, and each of the three functions on the right is
integrable: the first and second because $\varrho_m*g$, $\varrho_m*h$, $g$
and $h$ all lie in $L^{1}$ by Step 1, the third likewise.  Integrating and
using additivity of the integral, then (6) and (5), we get for every $m$
\begin{equation}\tag{7}
  \int\bigl|(\varrho_m*g)(x)-g(x)\bigr|\,d\sigma(x)
  \le\frac{\eta}{3}
    +\int\bigl|(\varrho_m*h)(x)-h(x)\bigr|\,d\sigma(x)
    +\frac{\eta}{3}.
\end{equation}
By Step 2 applied to this $h$, there is $M$ such that the middle term is
$<\eta/3$ for every $m\ge M$.  Hence
\[
  \int\bigl|(\varrho_m*g)(x)-g(x)\bigr|\,d\sigma(x)<\eta
  \qquad(m\ge M).
\]
Since $\eta>0$ was arbitrary and the quantity on the left is nonnegative, it
tends to $0$ as $m\to\infty$, which is the assertion.
\end{proof}
